\renewcommand{\bibname}{Quellen und Bandzuordnung}
\begin{thebibliography}{9}
\addcontentsline{toc}{chapter}{Quellen und Bandzuordnung}
\footnotesize
\setlength{\itemsep}{5pt}
\setlength{\parskip}{0pt}
\interlinepenalty=10000

\bibitem{project}
Martin Kunze: \emph{Die Grundlagen der Mathematik}.
Band~3: reine Mengenform des Auswahlaxioms
(\texttt{ChoiceFamilyAxiom}).
Band~9: Auswahlprinzip und seine Familienfassung.
Die Graphkonvention für Funktionen stammt aus Band~5.
Der vorliegende Leitfaden verwendet die Darstellungsform der
beweisfreien Hauptresultat-Leitfäden zu Band~3, Band~4, Band~5,
Band~28 und Band~37. Die Relations- und Surjektionsfassungen
des Hauptbands werden hier nicht entwickelt.

\bibitem{hrbacek}
Karel Hrbacek und Thomas Jech:
\emph{Introduction to Set Theory}. 3.~Auflage,
Marcel Dekker, 1999.
Kapitel~8, Abschnitt~1, Druckseiten~138--140;
insbesondere Satz~1.3 und die anschließende Produktform:
Auswahlfunktionen, Auswahlmengen disjunkter Familien,
indizierte Auswahl und Produkte.
Die dortige Konvention erlaubt beliebige Familien und
fordert die Auswahlbedingung nur für nichtleere Glieder;
hier werden leere Glieder bereits in der Voraussetzung ausgeschlossen.
\url{https://u.cs.biu.ac.il/~tsaban/Pdf/HJ_backup.pdf}.

\bibitem{spronk}
Nico Spronk: \emph{Axiom of Choice et al}.
Vorlesungsnotizen zu PMATH~453/753, University of Waterloo.
Seite~1: Auswahl auf allen nichtleeren Teilmengen einer Menge.
Seite~4: äquivalente Produktform und Abgrenzung zu endlichen Produkten.
Die in denselben Notizen behandelten Ordnungsprinzipien werden
für diesen Leitfaden nicht übernommen.
\url{https://www.math.uwaterloo.ca/~nspronk/math453/AofC.pdf}.

\bibitem{collins}
George E. Collins: \emph{Distributivity and an axiom of choice}.
The Journal of Symbolic Logic \textbf{19} (1954), Nr.~4,
S.~275--277.
Primärquelle zur Äquivalenz der vollständigen mengentheoretischen
Distributivität und ihrer dualen Fassung mit dem Auswahlprinzip.
\url{https://doi.org/10.2307/2267734}.

\end{thebibliography}

\noindent\textbf{Zuordnung und Formulierung der Aussagen.}
Kapitel~\ref{b09guide:ch:criteria} gibt die klassischen
Auswahlfassungen aus \cite{hrbacek,spronk} wieder.
Die Aussagen in Kapitel~\ref{b09guide:ch:selections}
sind mengenbasierte Anwendungsfassungen desselben Auswahlgrundsatzes.
Die beiden Schnitt- und Vereinigungsgesetze in
Kapitel~\ref{b09guide:ch:distribution} folgen der in
\cite{collins} behandelten vollständigen Distributivität;
das Produktgesetz ist ihre entsprechende Auswahlform für Tupel.

Die Notation mit variablen Indexmengen $J_i$, die expliziten
Randfälle und die Zusammenfassung mit vorgegebenen Werten
sind für diesen Leitfaden einheitlich formuliert.
Nicht jede Zusammenfassung steht in den Quellen wortgleich.
Es wird weder eine historische Erstentdeckung für jede
Anwendungsfassung noch Vollständigkeit aller bekannten
Auswahlfolgen beansprucht.
